%! TEX root = ./main.tex
% chktex-file 44:
% chktex-file 1
% chktex-file 36

\begin{Aufgabe}{3}
	\begin{tasks}
		\task Schreibe die Zahl \num{10010100001} als Wort.
		\task Schreibe fünfmillionendreihundertsechs als Dezimalzahl.
	\end{tasks}
\end{Aufgabe}
\begin{Loesung}
	\begin{tasks}
		\task zehnmilliardenzehnmillioneneinhunderttausendundeins \Punkt{}
		\task \num{500306} \Punkt
	\end{tasks}
\end{Loesung}
\begin{Aufgabe}{3}
	\begin{tasks}
		\task Runde die Zahl \num{453287}{} auf Hunderter, auf Tausender und auf Millionen.
		\task Melanie sagt stolz: "`Ich habe gestern beim Fußballspielen auf Zehner gerundet 20 Bälle gehalten!"'
		Gib an, wie viele Bälle sie mindestens und wie viele Bälle sie höchstens gehalten hat.
	\end{tasks}
\end{Aufgabe}
\begin{Loesung}
	\begin{tasks}(2)
		\task $
		\begin{aligned}
			\num{453287}
			&\approx \num{453300}\quad\text{auf Hunderter}\HalberPunkt{}\\
			&\approx \num{453000}\quad\text{auf Tausender}\HalberPunkt{}\\
			&\approx \num{0}\quad\text{auf Millionen}\Punkt{}\\
		\end{aligned}
		$
		\task Mindestens 15 \HalberPunkt{} und höchstens 24 \HalberPunkt{}
	\end{tasks}
\end{Loesung}

\begin{Aufgabe}{4}
	Führe erst eine Überschlagsrechnung durch. Berechne dann schriftlich.
	\begin{tasks}(3)
		\task $\num{4528}+\num{78634}$
		\task $\num{775423}-\num{534768}$
		\task $\num{14572}\cdot\num{37}$
	\end{tasks}
\end{Aufgabe}
\begin{Loesung}
	\begin{tasks}
		\task Ü: $\num{5000}+\num{79000}=\num{84000}$ \HalberPunkt;
		genau \num{83162} \HalberPunkt{}
		\task Ü: $\num{780000}-\num{530000}=\num{250000}$ \HalberPunkt;
		genau \num{240655} \Punkt{}
		\task Ü: $\num{15000}\cdot\num{40}=\num{600000}$ \HalberPunkt;
		genau \num{539164} \Punkt{}
	\end{tasks}
\end{Loesung}
\begin{Aufgabe}{3}
	Schreibe den Rechenterm auf und berechne das Ergebnis.
	\begin{tasks}
		\task Das Produkt der Zahlen 8 und 17.
		\task Addiere zu der Differenz der Zahlen 12 und 3 die Zahl 23.
	\end{tasks}
\end{Aufgabe}
\begin{Loesung}
	\begin{tasks}(2)
		\task $8\cdot 17 =136$\Punkt \task $12-3=9\Punkt\quad 9+23=32$\Punkt
	\end{tasks}
\end{Loesung}
\begin{Bonusaufgabe}{1}
	Wenn man ein Streichholz umlegt, wird die Rechnung richtig. Schreibe die richtige Rechnung auf.
	\begin{center}
	\begin{tikzpicture}[very thick]
		\draw (0.5, 0) -- (0.5, 2);

		\draw (1, 0) -- (2, 2);
		\draw (1, 2) -- (2, 0);

		\draw (2.5, 1) -- (4.5, 1);

		\draw (5, 2) -- (5.5, 0);
		\draw (6, 2) -- (5.5, 0);

		\draw (6.5, 0) -- (6.5, 2);

		\draw (7, 0.8) -- (9, 0.8);
		\draw (7, 1.2) -- (9, 1.2);

		\draw (9.5, 0) -- (10.5, 2);
		\draw (9.5, 2) -- (10.5, 0);

		\draw (11, 0) -- (11, 2);

		\draw (11.5, 2) -- (12, 0);
		\draw (12, 0) -- (12.5, 2);
	\end{tikzpicture}
	\end{center}
\end{Bonusaufgabe}
\begin{Loesung}
	\Rom{9}+\Rom{6}=\Rom{15}\Punkt
\end{Loesung}
